% ============================================================================
\section{Beyond Transformers: Alternative Architectures}
\label{sec:alt-architectures}
% ============================================================================

The previous section built a detailed understanding of the transformer
architecture and its attention mechanism. Transformers dominate modern
AI, but they carry a fundamental cost: the self-attention mechanism
scales \emph{quadratically} with sequence length. Processing a
sequence of $n$ tokens requires $O(n^2)$ operations and $O(n^2)$
memory for the attention matrix. At inference time, the
\concept{KV cache}---the stored keys and values from all previous
tokens---grows linearly with context length and must be held in fast
memory.

This quadratic bottleneck has motivated a wave of alternative
architectures that aim to match transformer quality while achieving
linear or even constant per-token cost. For our bitwise AI project,
these alternatives are especially interesting: many replace the
softmax attention (the hardest component to binarise, as discussed in
\cref{sec:transformers-llms}) with simple recurrent updates that
naturally map to integer arithmetic.

This section surveys the most promising alternatives in order of
relevance to bitwise computation, then compares them systematically.

% ----------------------------------------------------------------------------
\subsection{State Space Models: Mamba}
\label{sec:mamba}
% ----------------------------------------------------------------------------

\subsubsection{The Core Idea: Attention Without a Matrix}

Recall from \cref{sec:transformers-llms} that self-attention works by
having every token ``look at'' every other token through the $QK^T$
matrix. This is powerful but expensive: if you double the sequence
length, the attention computation quadruples.

\concept{State Space Models} (SSMs) take a fundamentally different
approach. Instead of computing pairwise relationships between all
tokens, an SSM maintains a compact \emph{hidden state} that
summarises everything seen so far. Each new token updates this state
through a simple recurrence:

\begin{equation}
  \mathbf{h}_t = \mathbf{A} \mathbf{h}_{t-1} + \mathbf{B} x_t,
  \qquad
  y_t = \mathbf{C} \mathbf{h}_t
  \label{eq:ssm-recurrence}
\end{equation}

where $\mathbf{h}_t$ is the hidden state at step $t$, $x_t$ is the
current input, and $y_t$ is the output. The matrices $\mathbf{A}$,
$\mathbf{B}$, and $\mathbf{C}$ control how the state evolves, how
input enters, and how output is read.

\begin{analogy}
Think of a transformer as a conference room where everyone can hear
everyone else---powerful, but the noise scales with the number of
people. An SSM is more like a relay runner passing a baton: each
runner (token) receives the baton (hidden state), adds their
contribution, and passes it forward. The baton has a fixed size
regardless of how many runners came before.
\end{analogy}

The key innovation of \concept{Mamba}~\cite{gu2024mamba} is making
the parameters $\mathbf{B}$, $\mathbf{C}$, and a discretisation step
$\boldsymbol{\Delta}$ \emph{input-dependent}---they change based on
what the current token is. This \concept{selective} mechanism allows
the model to decide what to remember and what to forget, analogous to
the gating in LSTMs but with much simpler arithmetic:

\begin{equation}
  \mathbf{B}_t = \text{Linear}(x_t), \quad
  \mathbf{C}_t = \text{Linear}(x_t), \quad
  \boldsymbol{\Delta}_t = \text{softplus}(\text{Linear}(x_t))
  \label{eq:mamba-selective}
\end{equation}

The discretised state transition becomes:
\begin{equation}
  \bar{\mathbf{A}}_t = \exp(\boldsymbol{\Delta}_t \mathbf{A}), \qquad
  \bar{\mathbf{B}}_t = \boldsymbol{\Delta}_t \mathbf{B}_t
  \label{eq:mamba-discretise}
\end{equation}

\begin{keyinsight}
The Mamba recurrence $\mathbf{h}_t = \bar{\mathbf{A}}_t
\mathbf{h}_{t-1} + \bar{\mathbf{B}}_t x_t$ is structurally a
\textbf{multiply-accumulate with decay}---the same operation as a
digital IIR filter. There is no softmax, no exponentiation over
sequences, no quadratic matrix. The entire inference step is a
matrix-vector multiply plus an addition. This is exactly the kind of
operation that maps naturally to integer and bitwise arithmetic.
\end{keyinsight}

\subsubsection{Architecture of a Mamba Block}

Each Mamba block follows a specific pattern, illustrated in
\cref{fig:mamba-block}:

\begin{figure}[htbp]
\centering
\begin{tikzpicture}[>=Stealth,
  comp/.style={draw, rounded corners=3pt, minimum height=0.7cm,
    minimum width=3.2cm, font=\small\sffamily, align=center,
    inner sep=4pt},
  arr/.style={-{Stealth[length=4pt]}, thick},
  plus/.style={circle, draw, minimum size=0.35cm, font=\tiny,
    fill=bitgreen!10, inner sep=0pt},
]
  % Input
  \node[comp, fill=bitgray!10] (input) at (0, 0) {Input $x$};

  % Two branches
  \node[comp, fill=bitblue!10] (lin1) at (-2, 1.5) {Linear proj.};
  \node[comp, fill=bitblue!10] (lin2) at (2, 1.5) {Linear proj.};

  % Left branch
  \node[comp, fill=bitorange!10] (conv) at (-2, 3) {Conv1d};
  \node[comp, fill=bitorange!10] (silu1) at (-2, 4.5) {SiLU};
  \node[comp, fill=bitpurple!15] (ssm) at (-2, 6)
    {SSM\\[-2pt]\scriptsize$(A, B_t, C_t, \Delta_t)$};

  % Right branch
  \node[comp, fill=bitorange!10] (silu2) at (2, 4.5) {SiLU};

  % Multiply gate
  \node[font=\Large] (mult) at (0, 6) {$\otimes$};

  % Output projection
  \node[comp, fill=bitblue!10] (outproj) at (0, 7.5) {Linear proj.};
  \node[comp, fill=bitgray!10] (output) at (0, 9) {Output};

  % Arrows
  \draw[arr] (input) -- (-2, 0.8) -- (lin1);
  \draw[arr] (input) -- (2, 0.8) -- (lin2);
  \draw[arr] (lin1) -- (conv);
  \draw[arr] (conv) -- (silu1);
  \draw[arr] (silu1) -- (ssm);
  \draw[arr] (ssm) -- (mult);
  \draw[arr] (lin2) -- (silu2);
  \draw[arr] (silu2) -- (2, 6) -- (mult);
  \draw[arr] (mult) -- (outproj);
  \draw[arr] (outproj) -- (output);

  % Annotations
  \node[font=\scriptsize\sffamily, bitblue, anchor=east] at (-3.8, 3)
    {local context};
  \node[font=\scriptsize\sffamily, bitpurple, anchor=east] at (-3.8, 6)
    {selective state space};
  \node[font=\scriptsize\sffamily, bitorange, anchor=west] at (3.8, 4.5)
    {gating branch};
\end{tikzpicture}
\caption[Architecture of a single Mamba block.]{%
  Architecture of a single Mamba block. The input is projected along
  two branches. The left branch applies a 1D convolution (for local
  context), a SiLU activation, and the selective SSM. The right branch
  provides a gating signal via SiLU. The element-wise product
  ($\otimes$) of the two branches is projected to produce the block
  output. Compare with the transformer block in
  \cref{fig:full-llm}---the SSM replaces the entire multi-head
  attention sub-layer.}
\label{fig:mamba-block}
\end{figure}

The block has two branches: a \emph{processing branch} (convolution
followed by the selective SSM) and a \emph{gating branch} (a simple
nonlinearity). These are combined by element-wise multiplication,
mirroring the gated linear unit (GLU) pattern common in modern
transformers.

\subsubsection{Complexity Comparison}

The computational advantage is stark:

\begin{center}
\begin{tabular}{lcc}
\toprule
\textbf{Operation} & \textbf{Transformer} & \textbf{Mamba} \\
\midrule
Training (per layer) & $O(n^2 d)$ & $O(n d s)$ \\
Inference per token & $O(n d)$ & $O(d s)$ \\
Memory (KV cache) & $O(n d)$ & $O(d s)$ \\
\bottomrule
\end{tabular}
\end{center}

\noindent where $n$ is sequence length, $d$ is model dimension, and
$s$ is the SSM state size (typically 16--64, independent of $n$).
At inference time, Mamba processes each token in $O(1)$ with respect
to sequence length---there is no KV cache that grows with context.

\subsubsection{Mamba-2, Mamba-3, and Scaling}

\concept{Mamba-2}~\cite{dao2024mamba2} introduced the
\concept{State Space Duality} (SSD) framework, showing that the
selective SSM can be equivalently computed as a structured
(semiseparable) matrix multiplication. This enables efficient hardware
utilisation via tensor cores, achieving 2--8$\times$ speedup over
the original Mamba while maintaining identical expressiveness.

\concept{Mamba-3} (ICLR 2026) extends the framework further with
complex-valued state matrices and MIMO (multiple-input,
multiple-output) state transitions, improving the model's ability
to track multiple independent ``threads'' of information
simultaneously.

The largest pure Mamba models include \concept{Falcon
Mamba}~7B~\cite{zuo2024falcon} and \concept{Codestral Mamba}~7B
(Mistral). Falcon Mamba~7B outperforms Llama~3~8B on
ARC, TruthfulQA, and GSM8K benchmarks while requiring
significantly less inference memory (no KV cache).

\subsubsection{The Recall Problem}

\begin{engineernote}
Pure SSMs have a fundamental limitation: the hidden state has a fixed
size. Once information is compressed into this state, it cannot be
perfectly recovered. Tasks requiring \concept{associative recall}---``what
word appeared after `elephant' three paragraphs ago?''---are difficult
because the model must have retained that specific association in its
finite state vector. Transformers, by contrast, store every token's
key and value in the KV cache and can look up any of them at any
time. This is not a minor issue: it is the primary motivation for the
hybrid architectures discussed in \cref{sec:hybrid-archs}.
\end{engineernote}

\subsubsection{Bitwise Relevance}

The Mamba recurrence is exceptionally well-suited to
integer arithmetic:

\begin{itemize}[leftmargin=2em]
  \item \textbf{Bi-Mamba}~\cite{ma2024bimamba} applies 1-bit
    weight quantisation to Mamba, achieving performance matching
    GPTQ 8-bit quantised transformers---with $8\times$ less weight
    memory.
  \item \textbf{LightMamba}~\cite{lightmamba2025} implements Mamba
    on FPGA at 4-bit precision, demonstrating $6\times$ energy
    efficiency improvement over GPU implementations.
  \item The core recurrence $\mathbf{h}_t = \bar{\mathbf{A}}
    \mathbf{h}_{t-1} + \bar{\mathbf{B}} x_t$ requires only
    multiply-accumulate---no softmax, no exponentiation, no
    division. Each of these operations can be replaced by
    fixed-point multiplication or, with appropriate quantisation,
    by bit-shifts and additions.
\end{itemize}

% ----------------------------------------------------------------------------
\subsection{RWKV: A Trainable RNN at Transformer Scale}
\label{sec:rwkv}
% ----------------------------------------------------------------------------

While Mamba replaces attention with a state space model,
\concept{RWKV}~\cite{peng2023rwkv} achieves a similar goal by
rethinking the recurrent neural network (RNN). Classical RNNs suffer
from vanishing gradients and cannot be parallelised during training.
RWKV solves both problems with a mechanism called
\concept{linear attention with exponential decay}.

\subsubsection{The WKV Mechanism}

The core operation in RWKV is the ``WKV'' (weighted key-value)
computation. For each token position $t$:

\begin{equation}
  \text{wkv}_t = \frac{
    e^{u + k_t} v_t + \sum_{i=1}^{t-1} e^{-(t-i)w + k_i} v_i
  }{
    e^{u + k_t} + \sum_{i=1}^{t-1} e^{-(t-i)w + k_i}
  }
  \label{eq:rwkv-wkv}
\end{equation}

where $w$ is a learned decay rate, $u$ is a bonus for the current
token, and $k_i, v_i$ are key and value vectors. The exponential
decay $e^{-(t-i)w}$ means that older tokens contribute
exponentially less---a built-in ``forgetting'' mechanism.

\begin{analogy}
If a transformer is a conference room (everyone hears everyone) and
Mamba is a relay race (pass the baton), then RWKV is like a
conversation where recent statements are crisp and clear, but older
ones fade into memory. You can still recall them, but with
decreasing clarity. The decay rate $w$ controls how fast the ``fog''
sets in.
\end{analogy}

The crucial property is that \cref{eq:rwkv-wkv} can be computed
\emph{either} as a parallel scan (for training) \emph{or} as a
recurrence (for inference). During training, all positions are
processed simultaneously with $O(n)$ total work. During inference,
each new token updates a fixed-size state in $O(1)$ time and
constant memory.

\subsubsection{RWKV-7 ``Goose'' and the Generalised Delta Rule}

The latest version, \concept{RWKV-7}~\cite{peng2025rwkv7}, replaces
the simple exponential decay with a \concept{Generalised Delta Rule}
that allows the model to learn more complex forgetting patterns.
RWKV-7 at 2.9B parameters matches the quality of Qwen2.5-3B (a
transformer) trained on $3\times$ more tokens.

RWKV has achieved remarkable real-world deployment:

\begin{itemize}[leftmargin=2em]
  \item Deployed on \textbf{1.5~billion Windows devices} for
    Microsoft Copilot edge scenarios, where constant-memory inference
    is critical.
  \item Open-source models from 0.1B to 14B parameters, trained by
    the RWKV Foundation community.
  \item The RWKV-7 architecture is available at sizes competitive
    with the latest transformer-based models.
\end{itemize}

\subsubsection{Bitwise Relevance}

RWKV is arguably the \emph{most} bitwise-amenable of all
alternative architectures:

\begin{itemize}[leftmargin=2em]
  \item \textbf{3-bit quantisation} with less than 1\% accuracy
    loss has been demonstrated, pushing close to the
    binary/ternary regime.
  \item \textbf{INT8 NEON kernels} are available for ARM
    processors, enabling efficient mobile deployment with integer
    arithmetic only.
  \item The WKV computation is dominated by multiply-accumulate
    with exponential decay---the same structure as a digital
    low-pass filter, which can be implemented with fixed-point
    arithmetic or even bit-shift approximations.
  \item The $O(1)$ inference memory means the model never builds a
    large dynamic data structure (like a KV cache) that would
    complicate fixed-point implementation.
\end{itemize}

\begin{keyinsight}
RWKV combines three properties that are rare together: transformer-competitive
quality, constant-memory inference, and demonstrated extreme
quantisation. For bitwise AI, this makes RWKV a strong candidate as
a backbone architecture---potentially easier to binarise than either
transformers or Mamba.
\end{keyinsight}

% ----------------------------------------------------------------------------
\subsection{Hybrid Architectures: Best of Both Worlds}
\label{sec:hybrid-archs}
% ----------------------------------------------------------------------------

The recall limitation of pure SSMs and RNNs
(\cref{sec:mamba}) has led to a pragmatic solution:
\concept{hybrid architectures} that interleave efficient recurrent
layers with a small number of attention layers. The attention layers
handle the tasks that require precise recall; the recurrent layers
handle everything else at linear cost.

\subsubsection{Jamba (AI21 Labs)}

\concept{Jamba}~\cite{lieber2024jamba} was the first large-scale
hybrid, combining Mamba layers, attention layers, and
mixture-of-experts (MoE) in a single model:

\begin{itemize}[leftmargin=2em]
  \item \textbf{398B total parameters, 94B active} per token (MoE
    routing activates only a subset of experts).
  \item Interleaves Mamba and attention layers at a ratio of roughly
    7:1 (one attention layer every 8 layers).
  \item At 256K context length, requires only \textbf{4\,GB of KV
    cache memory}---compared to $\sim$128\,GB for a pure transformer
    of similar size.
  \item Matches or exceeds Mixtral 8$\times$7B on standard
    benchmarks.
\end{itemize}

\subsubsection{Falcon H1 and H1R (TII)}

The \concept{Falcon H1}~\cite{tii2025falconh1} series takes a
different approach: instead of alternating layers, it places SSM and
attention \emph{heads} in \textbf{parallel within the same mixer
block}. Some heads attend, others recur, and their outputs are
combined:

\begin{itemize}[leftmargin=2em]
  \item Falcon H1 0.5B--34B parameter range, with the 7B model
    beating transformer models of 32B--47B parameters on several
    benchmarks.
  \item \textbf{Falcon H1R}~\cite{tii2026falconh1r} extends this
    with a recurrent-only variant for edge deployment.
  \item The parallel design avoids the ``attention every $k$ layers''
    pattern and allows more flexible resource allocation.
\end{itemize}

\subsubsection{Bamba (IBM)}

\concept{Bamba}~\cite{ibm2025bamba} combines Mamba-2 layers with
global attention layers, optimised for throughput:

\begin{itemize}[leftmargin=2em]
  \item Achieves $2\times$ throughput compared to pure transformer
    baselines of similar quality.
  \item Incorporated into IBM's Granite~4.0 model family.
  \item Demonstrates that the hybrid pattern is becoming an industry
    standard, not just a research curiosity.
\end{itemize}

\subsubsection{Kimi Linear (Moonshot AI)}

\concept{Kimi Linear}~\cite{kimi2025linear} replaces softmax attention
with \concept{linear attention}, achieving:

\begin{itemize}[leftmargin=2em]
  \item 75\% reduction in KV cache memory.
  \item The first linear-attention model to \emph{beat} full
    softmax attention on standard benchmarks at comparable scale.
  \item Demonstrates that the attention mechanism itself can be
    linearised without quality loss, given sufficient scale and
    training.
\end{itemize}

\begin{engineernote}
The hybrid pattern is converging across the industry: IBM (Bamba/Granite),
AI21 (Jamba), TII (Falcon H1), and Moonshot (Kimi) have all
independently adopted architectures that mix efficient recurrent
layers with selective attention. This is not a coincidence---it
reflects a fundamental architectural insight that different layers in
a deep network serve different functions, and not all of them need
the full power (and cost) of attention.
\end{engineernote}

% ----------------------------------------------------------------------------
\subsection{Other Notable Approaches}
\label{sec:other-alt-archs}
% ----------------------------------------------------------------------------

Several other architectures deserve mention for their distinct
contributions, even if they are less directly relevant to bitwise
computation.

\subsubsection{RetNet (Microsoft)}

\concept{RetNet}~\cite{sun2023retnet} introduces \concept{multi-scale
retention}, a mechanism that supports three computation modes:
parallel (for training), recurrent (for $O(1)$ inference), and
chunkwise (a hybrid). The multi-scale design uses different decay
rates at different heads, allowing the model to attend to both
short-range and long-range patterns simultaneously.

\subsubsection{xLSTM (Hochreiter)}

\concept{xLSTM}~\cite{beck2024xlstm}, from the original inventor of
the LSTM, modernises the classic architecture with exponential gating
and a matrix-valued memory cell. Key results include 40\% lower power
consumption than comparably-sized transformers, making it relevant for
energy-constrained deployment---though its more complex gating
mechanism is harder to binarise than Mamba or RWKV.

\subsubsection{Titans (Google DeepMind)}

\concept{Titans}~\cite{behrouz2025titans} augments attention with a
learned long-term memory module that can store and retrieve
information across very long contexts. It beats GPT-4 on
needle-in-a-haystack retrieval tasks, addressing the recall
limitation of SSMs from a different angle---by adding explicit
memory rather than falling back to attention.

\subsubsection{Mercury (Inception Labs)}

\concept{Mercury}~\cite{inception2025mercury} takes a radically
different approach: instead of generating tokens one at a time
(autoregressive), it uses a \concept{diffusion} process to generate
multiple tokens in parallel. This achieves throughput of
approximately 1{,}000 tokens per second on consumer hardware---an
order of magnitude faster than autoregressive models of similar size.
While not directly related to bitwise computation, it demonstrates
that the autoregressive bottleneck itself may be optional.

% ----------------------------------------------------------------------------
\subsection{Architecture Comparison}
\label{sec:alt-arch-comparison}
% ----------------------------------------------------------------------------

\Cref{tab:arch-comparison} provides a systematic comparison of all
architectures discussed in this section alongside the transformer
baseline.

\begin{table}[htbp]
\centering
\caption[Comparison of sequence modelling architectures.]{%
  Comparison of sequence modelling architectures. Complexity is given
  per layer, where $n$ is sequence length, $d$ is model dimension,
  and $s$ is SSM state size. ``Quality vs.\ Transformer'' indicates
  performance at comparable model size. Bitwise amenability reflects
  the ease of replacing floating-point operations with integer or
  binary arithmetic.}
\label{tab:arch-comparison}
\small
\begin{tabularx}{\textwidth}{@{}l c c c c c c@{}}
\toprule
\textbf{Architecture} &
\textbf{\shortstack{Train\\complexity}} &
\textbf{\shortstack{Inference\\per token}} &
\textbf{\shortstack{KV/state\\memory}} &
\textbf{\shortstack{Largest\\model}} &
\textbf{\shortstack{Production\\deployed?}} &
\textbf{\shortstack{Bitwise\\amenability}} \\
\midrule
Transformer &
  $O(n^2 d)$ & $O(nd)$ & $O(nd)$ &
  $>$1T & Yes &
  \textcolor{bitorange}{Medium} \\
Mamba (SSM) &
  $O(nds)$ & $O(ds)$ & $O(ds)$ &
  7B & Yes &
  \textcolor{bitgreen}{\textbf{High}} \\
RWKV &
  $O(nd)$ & $O(d)$ & $O(d)$ &
  14B & Yes &
  \textcolor{bitgreen}{\textbf{High}} \\
Jamba (hybrid) &
  $O(nd)^*$ & $O(d)^*$ & $O(nd)^{\dagger}$ &
  398B & Yes &
  \textcolor{bitorange}{Medium} \\
Falcon H1 &
  $O(nd)^*$ & $O(d)^*$ & $O(nd)^{\dagger}$ &
  34B & Yes &
  \textcolor{bitorange}{Medium} \\
RetNet &
  $O(nd)$ & $O(d)$ & $O(d)$ &
  6.7B & No &
  \textcolor{bitorange}{Medium} \\
xLSTM &
  $O(nd)$ & $O(d)$ & $O(d^2)$ &
  7B & No &
  \textcolor{bitorange}{Medium} \\
Titans &
  $O(n^2 d)$ & $O(nd)$ & $O(nd)$ &
  --- & No &
  \textcolor{red!70!black}{Low} \\
Mercury &
  $O(n^2 d)$ & $O(d)^{\ddagger}$ & $O(nd)$ &
  --- & Yes &
  \textcolor{red!70!black}{Low} \\
\bottomrule
\multicolumn{7}{@{}l}{%
  \scriptsize $^*$Amortised over recurrent layers; attention layers
  remain $O(n^2d)$.
  $^{\dagger}$Reduced by $\sim\!75$--$97\%$ vs.\ pure transformer.
  $^{\ddagger}$Parallel generation amortises cost.
} \\
\end{tabularx}
\end{table}

% ----------------------------------------------------------------------------
\subsection{Implications for This Project}
\label{sec:alt-arch-implications}
% ----------------------------------------------------------------------------

The alternative architectures surveyed above share a crucial
property: they replace the softmax attention bottleneck with
\concept{simple recurrent updates}---multiply-accumulate with decay,
weighted sums with exponential forgetting, or gated state
transitions. These are structurally identical to the operations
found in \concept{digital signal processing} (DSP): IIR filters,
exponential moving averages, and leaky integrators.

\begin{keyinsight}
The connection between SSMs and DSP is not a loose analogy---it is a
mathematical identity. The Mamba recurrence
$h_t = \bar{A} h_{t-1} + \bar{B} x_t$ \emph{is} a first-order IIR
filter. The RWKV exponential decay \emph{is} an exponential moving
average. These are operations that DSP engineers have implemented in
fixed-point integer arithmetic for decades, on hardware as simple as
8-bit microcontrollers. The ``bitwise bottleneck'' that makes
transformer attention so hard to binarise simply does not exist in
these architectures.
\end{keyinsight}

This observation connects directly to the sound synthesis vision
outlined in the project roadmap. Classical physical modelling synthesis
(Karplus-Strong, waveguide models) already operates with
near-bitwise operations: delay lines, low-pass filters (which are
just multiply-accumulate with decay), and feedback loops. An
SSM-based architecture for audio generation would use the
\emph{same} computational primitives as the synthesis algorithms it
aims to learn.

\Cref{fig:bitwise-ssm-mapping} illustrates how the four bitwise
paradigms from this article map onto SSM components---analogous to
the transformer mapping in \cref{sec:scaling-llms}.

\begin{figure}[htbp]
\centering
\begin{tikzpicture}[>=Stealth,
  comp/.style={draw, rounded corners=4pt, minimum height=1cm,
    minimum width=3.5cm, font=\small\sffamily, align=center,
    inner sep=5pt},
  paradigm/.style={draw, rounded corners=2pt, minimum height=0.6cm,
    font=\scriptsize\sffamily, align=center, inner sep=3pt,
    dashed, thick},
  arr/.style={-{Stealth[length=4pt]}, thick, bitgray},
  maparr/.style={-{Stealth[length=3pt]}, thick, dashed},
]
  % SSM components (left column)
  \node[comp, fill=bitblue!10] (embed) at (0, 8) {Embedding};
  \node[comp, fill=bitorange!10] (conv) at (0, 6.2) {Conv1d};
  \node[comp, fill=bitpurple!15] (ssm) at (0, 4.4)
    {SSM recurrence\\[-2pt]\scriptsize$h_t = Ah_{t-1} + Bx_t$};
  \node[comp, fill=bitgreen!10] (gate) at (0, 2.6) {Gating (SiLU)};
  \node[comp, fill=bitblue!10] (proj) at (0, 0.8) {Output projection};

  % Arrows between components
  \draw[arr] (embed) -- (conv);
  \draw[arr] (conv) -- (ssm);
  \draw[arr] (ssm) -- (gate);
  \draw[arr] (gate) -- (proj);

  % Bitwise paradigms (right column)
  \node[paradigm, fill=bitblue!8, text=bitblue!80!black]
    (bnn) at (7, 7.1) {BNN / BitNet\\ternary weights};
  \node[paradigm, fill=bitorange!8, text=bitorange!80!black]
    (lut) at (7, 5.3) {LUT compilation\\activation tables};
  \node[paradigm, fill=bitpurple!8, text=bitpurple!80!black]
    (lgn) at (7, 3.5) {Logic Gate Networks\\learned gate topology};
  \node[paradigm, fill=bitgreen!8, text=bitgreen!80!black]
    (tm) at (7, 1.7) {Tsetlin Machines\\clause-based gating};

  % Mapping arrows
  \draw[maparr, bitblue] (embed.east) -- ++(0.5,0) |- (bnn.west);
  \draw[maparr, bitblue] (proj.east) -- ++(0.3,0) |- (bnn.west);
  \draw[maparr, bitorange] (gate.east) -- ++(0.5,0) |- (lut.west);
  \draw[maparr, bitorange] (conv.east) -- ++(0.3,0) |- (lut.west);
  \draw[maparr, bitpurple] (ssm.east) -- (lgn.west);
  \draw[maparr, bitgreen] (gate.east) -- ++(0.8,0) |- (tm.west);

  % Title
  \node[font=\small\sffamily\bfseries, anchor=south] at (0, 8.8)
    {SSM Components};
  \node[font=\small\sffamily\bfseries, anchor=south] at (7, 8.0)
    {Bitwise Paradigms};
\end{tikzpicture}
\caption[Mapping of bitwise paradigms to SSM components.]{%
  How the four bitwise paradigms map to SSM (Mamba/RWKV) components.
  \textbf{BNN/BitNet} ternary weights apply to embedding and output
  projections. \textbf{LUT compilation} replaces activation functions
  (SiLU, softplus) and the Conv1d with lookup tables.
  \textbf{Logic Gate Networks} can implement the SSM recurrence
  itself as learned gate topology. \textbf{Tsetlin Machines} provide
  interpretable clause-based gating. Compare with the transformer
  mapping in \cref{sec:scaling-llms}.}
\label{fig:bitwise-ssm-mapping}
\end{figure}

The project recommendation is clear: SSM-based architectures---particularly
RWKV and Mamba---should be considered alongside transformers as
target architectures for bitwise scaling. Their recurrent structure
eliminates the softmax bottleneck entirely, their inference is
constant-memory (critical for edge deployment), and their core
operations are already proven to work at extreme quantisation levels.

% ----------------------------------------------------------------------------
% Full-page landscape diagram: Mamba architecture end-to-end
% ----------------------------------------------------------------------------

\begin{landscape}
\begin{figure}[p]
\centering
\begin{tikzpicture}[>=Stealth, x=0.84cm, y=0.84cm,
  % --- Shared styles ---
  comp/.style={draw, rounded corners=3pt, minimum height=0.65cm,
    font=\scriptsize\sffamily, align=center, inner sep=3pt},
  smallcomp/.style={draw, rounded corners=2pt, minimum height=0.5cm,
    font=\tiny\sffamily, align=center, inner sep=2pt},
  arr/.style={-{Stealth[length=4pt]}, thick},
  fatarr/.style={-{Stealth[length=5pt]}, very thick},
  plus/.style={circle, draw, minimum size=0.35cm, font=\tiny,
    fill=bitgreen!10, inner sep=0pt},
  brace/.style={decorate, decoration={brace, amplitude=5pt, mirror}},
  bracel/.style={decorate, decoration={brace, amplitude=5pt}},
  lbl/.style={font=\tiny\sffamily},
]

% ============================================================================
% COLUMN 1: Tokeniser + Embedding (same as transformer)
% ============================================================================
\node[font=\small\sffamily\bfseries, bitblue!80!black] at (2, 14.0)
  {Tokeniser + Embedding};

% Token boxes
\foreach \i/\tok in {0/The, 1/\_cat, 2/\_sat, 3/\_on, 4/\_the, 5/\_mat} {
  \node[comp, fill=bitgray!10, minimum width=1.1cm]
    (tok\i) at ({\i*0.85+0.2}, 12.8) {\tok};
}

% ID boxes
\foreach \i/\id in {0/464, 1/5023, 2/3156, 3/322, 4/279, 5/8643} {
  \node[smallcomp, fill=bitblue!8, minimum width=0.95cm]
    (id\i) at ({\i*0.85+0.2}, 11.9) {\id};
  \draw[arr, bitgray] (tok\i) -- (id\i);
}

% Embedding table
\node[comp, fill=bitblue!15, minimum width=5.0cm, minimum height=1.0cm]
  (emb) at (2.35, 10.3)
  {Embedding Table $\bm{E}$\\$V \times d$};
\foreach \i in {0,...,5} {
  \draw[arr, bitblue!60] (id\i) -- ({\i*0.85+0.2}, 10.85);
}

% Embedding output vectors
\node[comp, fill=bitblue!10, minimum width=5.0cm]
  (embout) at (2.35, 8.8)
  {$n$ vectors, $d$ dims each};
\draw[arr] (emb) -- (embout);

% ============================================================================
% COLUMN 2: Mamba Block (expanded)
% ============================================================================
\node[font=\small\sffamily\bfseries, bitorange!80!black] at (12.5, 14.0)
  {Mamba Block (expanded)};

% Input
\node[comp, fill=bitgray!10, minimum width=2.4cm]
  (blkin) at (12.5, 12.8) {Block input $x$};

% Layer norm
\node[comp, fill=bitgray!15, minimum width=2.4cm]
  (norm1) at (12.5, 11.6) {RMSNorm};
\draw[arr] (blkin) -- (norm1);

% Two linear projections (branches)
\node[comp, fill=bitblue!12, minimum width=2.1cm]
  (linA) at (10.5, 10.3) {Linear ($d \!\to\! Ed$)};
\node[comp, fill=bitblue!12, minimum width=2.1cm]
  (linB) at (14.5, 10.3) {Linear ($d \!\to\! Ed$)};
\draw[arr] (norm1.south) -- ++(0,-0.25) -| (linA.north);
\draw[arr] (norm1.south) -- ++(0,-0.25) -| (linB.north);

% Left branch: Conv1d
\node[comp, fill=bitorange!12, minimum width=2.1cm]
  (conv1d) at (10.5, 8.9) {Conv1d ($k\!=\!4$)};
\draw[arr] (linA) -- (conv1d);

% Left branch: SiLU
\node[comp, fill=bitorange!12, minimum width=2.1cm]
  (silu1) at (10.5, 7.7) {SiLU $\sigma(x) \!\cdot\! x$};
\draw[arr] (conv1d) -- (silu1);

% SSM parameters
\node[comp, fill=bitpurple!15, minimum width=2.1cm, minimum height=0.9cm]
  (ssmparams) at (10.5, 6.2)
  {Selection\\[-2pt]\tiny $B_t, C_t, \Delta_t = f(x_t)$};
\draw[arr] (silu1) -- (ssmparams);

% SSM core
\node[comp, fill=bitpurple!20, minimum width=2.1cm, minimum height=1.1cm]
  (ssmcore) at (10.5, 4.4)
  {SSM scan\\[-2pt]\tiny $h_t = \bar{A}_t h_{t\text{-}1} + \bar{B}_t x_t$\\[-2pt]
   \tiny $y_t = C_t h_t$};
\draw[arr] (ssmparams) -- (ssmcore);

% A parameter (fixed, diagonal)
\node[smallcomp, fill=bitpurple!8, minimum width=1.3cm]
  (Aparam) at (8, 5.3) {\tiny $A$ (learned,\\[-2pt]\tiny diagonal)};
\draw[arr, bitpurple!60] (Aparam) -- (ssmcore);

% Right branch: SiLU (gating)
\node[comp, fill=bitorange!12, minimum width=2.1cm]
  (silu2) at (14.5, 7.7) {SiLU (gate)};
\draw[arr] (linB) -- ++(0,-1.9) -- (silu2);

% Element-wise multiply
\node[font=\normalsize, inner sep=2pt] (mult) at (12.5, 3.1) {$\otimes$};
\draw[arr] (ssmcore) -- ++(0,-0.45) -| (mult);
\draw[arr] (silu2) -- ++(0,-3.95) -| (mult);

% Output projection
\node[comp, fill=bitblue!12, minimum width=2.4cm]
  (outproj) at (12.5, 2.0) {Linear ($Ed \!\to\! d$)};
\draw[arr] (mult) -- (outproj);

% Residual connection
\node[plus] (res1) at (12.5, 1.0) {$+$};
\draw[arr] (outproj) -- (res1);
\draw[arr, dashed, bitgreen!60!black, rounded corners=5pt]
  (blkin.east) -- ++(1.5, 0) -- ++(0, -11.8) -- (res1.east);

% Block output
\node[comp, fill=bitgray!10, minimum width=2.4cm]
  (blkout) at (12.5, 0.0) {Block output};
\draw[arr] (res1) -- (blkout);

% Annotations
\node[lbl, bitorange, anchor=east, text width=1.5cm, align=right]
  at (8.5, 8.9) {local\\context};
\node[lbl, bitpurple, anchor=east, text width=2cm, align=right]
  at (8.0, 6.2) {input-\\dependent};
\node[lbl, bitpurple, anchor=east, text width=2cm, align=right]
  at (8.0, 4.4) {selective\\scan};
\node[lbl, bitorange, anchor=west, text width=2cm, align=left] at (16, 7.7)
  {gating\\branch};
\node[lbl, bitgreen!60!black, anchor=west] at (14.5, 0.7)
  {residual};

% ============================================================================
% COLUMN 3: Stacked blocks + Output head
% ============================================================================
\node[font=\small\sffamily\bfseries, bitgreen!70!black] at (22, 14.0)
  {Stack + Output Head};

% Stacked blocks
\node[comp, fill=bitorange!8, minimum width=3.4cm, minimum height=0.65cm]
  (b1) at (22, 12.5) {Mamba Block 1};
\node[comp, fill=bitorange!8, minimum width=3.4cm, minimum height=0.65cm]
  (b2) at (22, 11.4) {Mamba Block 2};
\node[font=\large, bitgray] at (22, 10.5) {$\vdots$};
\node[comp, fill=bitorange!8, minimum width=3.4cm, minimum height=0.65cm]
  (bN) at (22, 9.6) {Mamba Block $N$};

\draw[arr] (b1) -- (b2);
\draw[arr, bitgray, densely dotted] (b2) -- ++(0, -0.55);
\draw[arr, bitgray, densely dotted] (22, 10.2) -- (bN);

% Final norm
\node[comp, fill=bitgray!15, minimum width=3.4cm]
  (finalnorm) at (22, 8.3) {RMSNorm};
\draw[arr] (bN) -- (finalnorm);

% Output projection
\node[comp, fill=bitblue!12, minimum width=3.4cm]
  (vocabproj) at (22, 7.1) {$\bm{W}_{\text{vocab}}$ ($d \!\to\! V$)};
\draw[arr] (finalnorm) -- (vocabproj);

% Logits
\node[comp, fill=bitgray!8, minimum width=3.4cm]
  (logits) at (22, 5.9) {Logits (size $V$)};
\draw[arr] (vocabproj) -- (logits);

% Softmax
\node[comp, fill=bitgreen!12, minimum width=3.4cm]
  (softmax) at (22, 4.7) {Softmax};
\draw[arr] (logits) -- (softmax);

% Probabilities
\node[comp, fill=bitgreen!8, minimum width=3.4cm]
  (probs) at (22, 3.5) {Probabilities};
\draw[arr] (softmax) -- (probs);

% Sample
\node[comp, fill=bitgreen!15, minimum width=3.4cm]
  (sample) at (22, 2.3) {Sample next token};
\draw[arr] (probs) -- (sample);

% ============================================================================
% CONNECTING ARROWS between columns
% ============================================================================

% Embedding to Block 1
\draw[fatarr] (embout.east) -- ++(1, 0) |- (b1.west);

% Block detail connection (from column 2 to column 3 stack)
% Route along right side of column 2, avoiding low y-values
\draw[fatarr, bitorange!50, dashed]
  (blkout.east) -- ++(4, 0) -- ++(0, 12.5) -- (b1.south east);
\node[lbl, bitorange, fill=white, inner sep=1pt] at (17.5, 6.5)
  {detail of each block};

% Autoregressive loop
\draw[fatarr, bitgreen!70!black, rounded corners=8pt]
  (sample.south) -- ++(0, -0.8) -- (2, 0.8) -- (2, 8.2);
\node[lbl, bitgreen!70!black, fill=white, inner sep=2pt]
  at (12, 0.8) {autoregressive loop: append token, repeat};

% ============================================================================
% TITLE
% ============================================================================
\node[font=\small\sffamily\bfseries, bitblue!80!black, anchor=north west]
  at (0, 14.8)
  {The Complete Mamba Architecture: End-to-End Pipeline};

% Dimension annotations
\node[lbl, bitgray, text width=3cm, align=center] at (2.35, 9.5)
  {$n$ tokens\\$d$ dims each};

% State size annotation
\node[lbl, bitpurple, text width=2.5cm, align=center] at (8, 3.7)
  {state size $s$\\(typically 16--64)};

\end{tikzpicture}
\caption[The complete Mamba architecture in a single view.]{%
  The complete Mamba architecture in a single view. \textbf{Left:}
  the tokeniser and embedding layer are identical to the transformer
  (\cref{fig:llm-full-picture}). \textbf{Centre:} a single Mamba
  block is expanded, showing the dual-branch structure: the
  processing branch (Conv1d for local context, input-dependent
  selection of $B_t$, $C_t$, $\Delta_t$, and the selective SSM scan)
  and the gating branch (SiLU activation). The element-wise product
  $\otimes$ merges both branches before the output projection. A
  residual connection bypasses the entire block. \textbf{Right:}
  $N$ blocks are stacked, followed by final RMSNorm and the output
  head (vocabulary projection, softmax, sampling). The autoregressive
  loop feeds each generated token back as input. Note the absence of
  any $O(n^2)$ attention matrix---the entire sequence-mixing
  computation is the $O(ds)$ SSM scan.}
\label{fig:mamba-full-picture}
\end{figure}
\end{landscape}
